\scanpage{323}
% archive sheet number: 324

$(1)\Longleftrightarrow(2)$ déjà vu.

$(2)\Longleftrightarrow(3)$ La donnée d'\unsure{un hom de|hom hom de}
foncteurs cov
\[
F_A\xrightarrow{u_p}\deleted{\unsure{F_p}}G=F_A\circ\Gamma^p
\]
équiv. à la donnée d'un élém
\[
u_p\in\deleted{\unsure{F_A}}G(A)=\operatorname{Hom}(A,\Gamma^{p}(A))
\]
Il faut exprimer les cond. de (2).

Les \unsure{Cond'ns} $f^{(0)}=1$, $f^{(1)}=\unsure{0|o}$ expriment
$u_0^{\deleted{\unsure{(0)}}}=1$, $u_1^{\deleted{\unsure{(1)}}}=\mathrm{id}$.

La relation $(f+g)^{(p)}=\sum f^{(i)}g^{(j)}$ : il
suffit de voir le cas universel
\[
\begin{aligned}
M&=A\times A,\quad f=\mathrm{inj}_{1}=\unsure{\alpha},\quad
g=\mathrm{inj}_{2}\deleted{=\unsure{\beta}},\\
(f+g)\deleted{{}^{\unsure{p}}}&=\Delta:A\longrightarrow A\times A=M
\end{aligned}
\]
\[
(f+g)^{\unsure{p|(p)}}=\Gamma^{p}(\Delta)\circ u_p\quad
\deleted{\Gamma^{p}(A)\longrightarrow}\Gamma^{p}(M)
\]

\diagram{\small\begin{tikzcd}[ampersand replacement=\&,column sep=large,row sep=large]
A \arrow[r,"{u_p}"]\arrow[d,"{\Delta_A^{(p)}}"] \&
\Gamma^{p}(A)\arrow[r,"{\Gamma^{p}(\Delta_{\unsure{A}})}"]
\arrow[dr,"{\Delta_{\Gamma^{\ast}A}}"'] \&
\Gamma^{p}(A\times\unsure{M|A})\arrow[d,"\simeq"]\\
\underbrace{A\times\cdots\times A}_{p} \arrow[rr,"{\text{car } A\longrightarrow\Gamma^{i}(A)\otimes\Gamma^{\unsure{j)|j}}(A)}"'] \&
{} \& \displaystyle\sum_{i+j=p}\Gamma^{\unsure{i|pi}}A\otimes\Gamma^{j}A
\end{tikzcd}}

On trouve $\Delta(u_p)=\displaystyle\sum_{i+j=p}u_i\otimes u_j$
